\documentclass[12pt,a4paper]{article}

\usepackage[utf8]{inputenc}
\usepackage[T1]{fontenc}
\usepackage[french]{babel}
\usepackage{geometry}
\usepackage{xcolor}
\usepackage{longtable}
\usepackage{array}
\usepackage{tabularx}
\usepackage{enumitem}
\usepackage{titlesec}
\usepackage{fancyhdr}
\usepackage{colortbl}
\usepackage{hyperref}
\usepackage{listings}
\usepackage{pifont}
\usepackage{longtable}
\usepackage{colortbl}
\usepackage{xcolor}
\usepackage{ltablex}
\keepXColumns


\geometry{margin=2.2cm}

% =========================
% CHARTE
% =========================
\definecolor{EspritBleu}{RGB}{0,102,153}
\definecolor{GrisClair}{RGB}{240,240,240}
\definecolor{GrisCode}{RGB}{245,245,245}

\newcommand{\NomEcole}{ESPRIT School of Business}
\newcommand{\Departement}{Département Informatique et Mathématiques Appliquées}
\newcommand{\cbox}{\ding{110}}

\pagestyle{fancy}
\fancyhf{}
\fancyhead[L]{\textbf{\NomEcole}}
\fancyhead[R]{\textbf{Architecture Multi-Agents IA}}
\fancyfoot[C]{\Departement}
\fancyfoot[R]{\thepage}

\titleformat{\section}
{\color{white}\normalfont\Large\bfseries}
{}{0em}
{}
[\color{EspritBleu}\titlerule]

\titleformat{\subsection}
{\normalfont\large\bfseries\color{EspritBleu}}
{\thesubsection}{1em}{}

\lstset{
    basicstyle=\ttfamily\small,
    backgroundcolor=\color{GrisCode},
    frame=single,
    breaklines=true,
    columns=fullflexible,
    showstringspaces=false
}

\begin{document}

\begin{center}
    {\LARGE \textbf{\NomEcole}}\\[0.2cm]
    {\large \Departement}\\[0.5cm]
    {\LARGE \textcolor{EspritBleu}{\textbf{Architecture Multi-Agents IA pour l’Analyse et la Refonte Automatique d’un Examen}}}\\[0.3cm]
    {\large Génération d’un rapport final et d’un examen restructuré à partir de fichiers \texttt{.tex}}
\end{center}

\vspace{0.5cm}

\section*{\colorbox{EspritBleu}{\parbox{\dimexpr\linewidth-2\fboxsep}{\centering \color{white}1. Objectif du Système}}}

Ce document décrit une architecture multi-agents permettant :
\begin{itemize}[leftmargin=1cm]
    \item d’analyser un examen réel au format PDF, DOCX, image ou texte,
    \item d’extraire les informations académiques, pédagogiques et techniques,
    \item de classifier les questions selon leur type et leur niveau cognitif,
    \item de réaliser le mapping entre les questions et les acquis d’apprentissage (AA),
    \item d’estimer la durée de traitement de chaque question,
    \item de proposer automatiquement des améliorations,
    \item de générer de nouvelles questions, de nouveaux exercices et de nouvelles activités pratiques,
    \item de réorganiser une nouvelle proposition d’examen en supprimant, remplaçant ou insérant des éléments,
    \item de générer le rapport officiel d’évaluation d’examen à partir d’un fichier \texttt{rapport examen officiel.tex},
    \item de générer la version finale de l’examen sur une plateforme cible,
    \item de notifier tous les enseignants du cours après génération ou mise à jour.
\end{itemize}

\section*{\colorbox{EspritBleu}{\parbox{\dimexpr\linewidth-2\fboxsep}{\centering \color{white}2. Vue d’Ensemble de l’Architecture}}}

L’architecture repose sur un \textbf{orchestrateur central} qui pilote une chaîne d’agents spécialisés.  
Chaque agent reçoit une entrée structurée, effectue une tâche précise, puis renvoie une sortie normalisée en JSON.  
Les agents d’analyse produisent une représentation détaillée de l’examen existant.  
Les agents de refonte produisent une proposition révisée.  
Les agents de rendu produisent le rapport officiel, le fichier final de l’examen et les exports vers la plateforme.

\subsection*{Chaîne de traitement globale}

\begin{enumerate}[leftmargin=1cm]
    \item Ingestion du document d’examen
    \item OCR / reconstruction si nécessaire
    \item Extraction des métadonnées
    \item Analyse de la forme de l’examen
    \item Segmentation des questions
    \item Classification automatique des questions
    \item Extraction du barème
    \item Mapping AA
    \item Analyse selon la taxonomie de Bloom
    \item Estimation de la durée par question
    \item Calcul des indicateurs qualité
    \item Analyse de difficulté et adéquation durée / charge
    \item Génération de propositions d’amélioration
    \item Génération de nouvelles questions, exercices et pratiques
    \item Réorganisation de la nouvelle proposition d’examen
    \item Vérification de conformité institutionnelle
    \item Consolidation des résultats
    \item Génération du rapport d’évaluation d’examen via \texttt{rapport examen officiel.tex}
    \item Génération LaTeX de sortie
    \item Validation et compilation LaTeX
    \item Génération de l’examen sur la plateforme cible
    \item Notification à tous les enseignants du cours
\end{enumerate}

\section*{\colorbox{EspritBleu}{\parbox{\dimexpr\linewidth-2\fboxsep}{\centering \color{white}3. Description Détaillée des Agents}}}

\subsection*{Agent 1 : Orchestrateur Central}

\textbf{Rôle :}
\begin{itemize}[leftmargin=1cm]
    \item piloter l’ensemble du workflow,
    \item appeler les agents dans l’ordre logique,
    \item gérer les relances, erreurs et validations intermédiaires,
    \item consolider les sorties finales.
\end{itemize}

\textbf{Entrées :}
\begin{itemize}[leftmargin=1cm]
    \item fichier de l’examen,
    \item fichier \texttt{template.tex},
    \item fichier \texttt{rapport examen officiel.tex},
    \item référentiel des acquis d’apprentissage,
    \item règles institutionnelles,
    \item métadonnées du cours,
    \item paramètres de plateforme cible.
\end{itemize}

\textbf{Sorties :}
\begin{itemize}[leftmargin=1cm]
    \item JSON consolidé,
    \item fichier \texttt{report\_final.tex},
    \item fichier \texttt{rapport\_evaluation\_final.tex},
    \item PDF compilé,
    \item structure finale de la nouvelle proposition d’examen,
    \item journal d’exécution.
\end{itemize}

\subsection*{Agent 2 : Agent d’Ingestion Documentaire}

\textbf{Rôle :}
\begin{itemize}[leftmargin=1cm]
    \item détecter le type de document,
    \item extraire le texte brut page par page,
    \item repérer tableaux, images, annexes et structures visibles.
\end{itemize}

\textbf{Formats pris en charge :}
\begin{itemize}[leftmargin=1cm]
    \item PDF natif,
    \item PDF scanné,
    \item image,
    \item DOCX,
    \item texte brut.
\end{itemize}

\subsection*{Agent 3 : Agent OCR / Reconstruction}

\textbf{Rôle :}
\begin{itemize}[leftmargin=1cm]
    \item reconstruire la structure si le document est scanné,
    \item détecter les blocs : entête, consignes, questions, sous-questions, tableaux, annexes.
\end{itemize}

\textbf{Sortie attendue :}
\begin{itemize}[leftmargin=1cm]
    \item structure logique des blocs du document,
    \item texte nettoyé et segmenté.
\end{itemize}

\subsection*{Agent 4 : Agent d’Extraction de Métadonnées}

\textbf{Rôle :}
\begin{itemize}[leftmargin=1cm]
    \item extraire le nom de l’épreuve,
    \item identifier la classe,
    \item détecter la langue,
    \item repérer la durée,
    \item repérer la date,
    \item identifier l’enseignant ou les enseignants,
    \item compter le nombre de pages.
\end{itemize}

\subsection*{Agent 5 : Agent d’Analyse de la Forme de l’Examen}

\textbf{Rôle :}
\begin{itemize}[leftmargin=1cm]
    \item détecter si l’étudiant répond sur la feuille,
    \item détecter si les documents sont autorisés,
    \item détecter si la calculatrice est autorisée,
    \item détecter si le PC est autorisé,
    \item détecter si Internet est autorisé,
    \item évaluer si la durée globale est adéquate,
    \item vérifier si tous les enseignants ont validé l’examen.
\end{itemize}

\subsection*{Agent 6 : Agent de Segmentation des Questions}

\textbf{Rôle :}
\begin{itemize}[leftmargin=1cm]
    \item identifier les exercices,
    \item séparer les questions principales et les sous-questions,
    \item donner un identifiant unique à chaque question ou exercice.
\end{itemize}

\subsection*{Agent 7 : Agent de Classification Automatique des Questions}

\textbf{Rôle :}
\begin{itemize}[leftmargin=1cm]
    \item attribuer un type à chaque élément évaluatif.
\end{itemize}

\textbf{Catégories possibles :}
\begin{itemize}[leftmargin=1cm]
    \item QCM,
    \item rédaction,
    \item pratique,
    \item étude de cas,
    \item exercice d’application,
    \item problème,
    \item projet,
    \item démonstration,
    \item calcul.
\end{itemize}

\subsection*{Agent 8 : Agent d’Extraction du Barème}

\textbf{Rôle :}
\begin{itemize}[leftmargin=1cm]
    \item détecter le nombre de points par question,
    \item détecter le nombre de points par sous-question,
    \item détecter le total global,
    \item vérifier si le barème détaillé est mentionné.
\end{itemize}

\subsection*{Agent 9 : Agent de Mapping AA}

\textbf{Rôle :}
\begin{itemize}[leftmargin=1cm]
    \item associer chaque question ou exercice aux acquis d’apprentissage concernés,
    \item vérifier la couverture globale des AA.
\end{itemize}

\subsection*{Agent 10 : Agent Bloom Taxonomy}

\textbf{Rôle :}
\begin{itemize}[leftmargin=1cm]
    \item classifier chaque question selon le niveau cognitif de Bloom.
\end{itemize}

\subsection*{Agent 11 : Agent d’Estimation de la Durée par Question}

\textbf{Rôle :}
\begin{itemize}[leftmargin=1cm]
    \item estimer le temps nécessaire pour traiter chaque question, exercice ou activité pratique,
    \item calculer une durée totale estimée de l’examen,
    \item signaler les questions trop longues ou trop courtes,
    \item alimenter l’analyse de l’adéquation entre durée annoncée et charge réelle.
\end{itemize}

\subsection*{Agent 12 : Agent d’Indicateurs Qualité}

\textbf{Rôle :}
\begin{itemize}[leftmargin=1cm]
    \item calculer les indicateurs qualité de l’examen.
\end{itemize}

\subsection*{Agent 13 : Agent d’Analyse de Difficulté}

\textbf{Rôle :}
\begin{itemize}[leftmargin=1cm]
    \item estimer la difficulté de chaque question,
    \item produire un niveau global de difficulté,
    \item croiser difficulté et durée estimée.
\end{itemize}

\subsection*{Agent 14 : Agent de Génération des Propositions d’Amélioration}

\textbf{Rôle :}
\begin{itemize}[leftmargin=1cm]
    \item proposer des améliorations ciblées sur l’examen existant,
    \item identifier les questions ambiguës, redondantes ou mal calibrées,
    \item suggérer des reformulations, suppressions, simplifications ou renforcements,
    \item produire une liste priorisée d’actions correctives.
\end{itemize}

\subsection*{Agent 15 : Agent de Génération de Questions, d’Exercices et de Pratiques}

\textbf{Rôle :}
\begin{itemize}[leftmargin=1cm]
    \item générer de nouveaux éléments évaluatifs adaptés aux besoins identifiés,
    \item produire des questions, des exercices complets, des études de cas et des activités pratiques,
    \item respecter les AA, le niveau Bloom visé, le barème et la durée cible.
\end{itemize}

\subsection*{Agent 16 : Agent de Réorganisation de la Nouvelle Proposition}

\textbf{Rôle :}
\begin{itemize}[leftmargin=1cm]
    \item construire une nouvelle version cohérente de l’examen,
    \item supprimer une ancienne question si nécessaire,
    \item insérer une nouvelle question, un nouvel exercice ou une nouvelle activité pratique,
    \item réordonner les éléments,
    \item préserver l’équilibre global en durée, difficulté et couverture pédagogique.
\end{itemize}

\subsection*{Agent 17 : Agent de Conformité Institutionnelle}

\textbf{Rôle :}
\begin{itemize}[leftmargin=1cm]
    \item vérifier la conformité à la charte de l’établissement,
    \item détecter les éléments manquants dans l’épreuve ou dans le rapport.
\end{itemize}

\subsection*{Agent 18 : Agent de Consolidation des Résultats}

\textbf{Rôle :}
\begin{itemize}[leftmargin=1cm]
    \item fusionner les résultats de tous les agents,
    \item produire un JSON final unique,
    \item résoudre les conflits éventuels entre agents.
\end{itemize}

\subsection*{Agent 19 : Agent Générateur du Rapport d’Évaluation d’Examen}

\textbf{Rôle :}
\begin{itemize}[leftmargin=1cm]
    \item générer le rapport officiel d’évaluation d’examen à partir du fichier \texttt{rapport examen officiel.tex},
    \item injecter les résultats consolidés dans ce modèle officiel,
    \item produire un fichier LaTeX dédié au rapport d’évaluation,
    \item séparer le rapport institutionnel d’évaluation de la génération technique de l’examen.
\end{itemize}

\textbf{Entrées :}
\begin{itemize}[leftmargin=1cm]
    \item fichier \texttt{rapport examen officiel.tex},
    \item JSON consolidé,
    \item indicateurs qualité,
    \item mapping AA,
    \item analyse Bloom,
    \item barème,
    \item plan de réorganisation.
\end{itemize}

\textbf{Sorties :}
\begin{itemize}[leftmargin=1cm]
    \item \texttt{rapport\_evaluation\_final.tex},
    \item \texttt{rapport\_evaluation\_final.pdf}.
\end{itemize}

\textbf{Variables injectées typiques :}
\begin{lstlisting}
{{ exam_header.exam_name }}
{{ exam_header.class }}
{{ exam_format.documents_allowed }}
{{ quality_indicators }}
{{ aa_mapping }}
{{ bloom_mapping }}
{{ improvement_proposals }}
\end{lstlisting}

\subsection*{Agent 20 : Agent de Génération LaTeX}

\textbf{Rôle :}
\begin{itemize}[leftmargin=1cm]
    \item injecter les données consolidées dans le fichier \texttt{template.tex},
    \item remplacer les marqueurs ou placeholders,
    \item générer un fichier \texttt{report\_final.tex}.
\end{itemize}

\subsection*{Agent 21 : Agent de Validation LaTeX}

\textbf{Rôle :}
\begin{itemize}[leftmargin=1cm]
    \item vérifier que les fichiers \texttt{.tex} générés compilent correctement,
    \item corriger les caractères spéciaux,
    \item sécuriser les champs sensibles : \_, \%, \&, \#,
    \item détecter les tableaux cassés ou les environnements manquants.
\end{itemize}

\subsection*{Agent 22 : Agent de Génération de l’Examen sur la Plateforme}

\textbf{Rôle :}
\begin{itemize}[leftmargin=1cm]
    \item transformer la structure finale de l’examen en format compatible avec la plateforme cible,
    \item créer les sections, questions, exercices et activités pratiques dans la plateforme,
    \item injecter les consignes, le barème, les durées et les métadonnées,
    \item préparer un export ou un déploiement automatisé.
\end{itemize}

\subsection*{Agent 23 : Agent de Notification aux Enseignants du Cours}

\textbf{Rôle :}
\begin{itemize}[leftmargin=1cm]
    \item notifier tous les enseignants rattachés au cours,
    \item transmettre l’état du traitement : analyse terminée, rapport généré, examen mis à jour, publication plateforme,
    \item envoyer les liens ou pièces jointes utiles,
    \item journaliser les notifications envoyées et leur statut.
\end{itemize}

\textbf{Entrées :}
\begin{itemize}[leftmargin=1cm]
    \item liste des enseignants du cours,
    \item métadonnées du cours,
    \item lien ou chemin du rapport d’évaluation,
    \item lien ou chemin de la nouvelle version de l’examen,
    \item statut de publication sur la plateforme.
\end{itemize}

\textbf{Sorties :}
\begin{itemize}[leftmargin=1cm]
    \item messages de notification,
    \item historique d’envoi,
    \item états : envoyé, en attente, erreur.
\end{itemize}

\textbf{Canaux possibles :}
\begin{itemize}[leftmargin=1cm]
    \item courrier électronique,
    \item notification plateforme,
    \item messagerie interne,
    \item tableau de bord enseignant.
\end{itemize}

\section*{\colorbox{EspritBleu}{\parbox{\dimexpr\linewidth-2\fboxsep}{\centering \color{white}4. Structure des Données Entre Agents}}}

Chaque agent doit renvoyer une structure standardisée.

\begin{lstlisting}
{
  "status": "ok",
  "confidence": 0.91,
  "data": {},
  "warnings": [],
  "errors": []
}
\end{lstlisting}

\subsection*{Exemple pour une question}

\begin{lstlisting}
{
  "status": "ok",
  "confidence": 0.88,
  "data": {
    "question_id": "Q2",
    "type": "Etude de cas",
    "points": 10,
    "bloom_level": "Analyze",
    "aa_ids": ["AA2", "AA3"],
    "estimated_time_minutes": 18
  },
  "warnings": ["Bareme detaille non explicite"],
  "errors": []
}
\end{lstlisting}

\section*{\colorbox{EspritBleu}{\parbox{\dimexpr\linewidth-2\fboxsep}{\centering \color{white}5. Exemple de JSON Consolidé}}}

\begin{lstlisting}
{
  "report_data": {
    "exam_header": {
      "exam_name": "Examen Final - Analyse",
      "class": "2INFO-B",
      "language": "Francais",
      "duration": "120 min",
      "date": "2026-01-18",
      "teachers": ["Dr X", "Mme Y"],
      "page_count": 5
    },
    "exam_format": {
      "answer_on_exam_sheet": true,
      "documents_allowed": false,
      "calculator_allowed": true,
      "pc_allowed": false,
      "internet_allowed": false,
      "duration_adequate": true
    },
    "questions": [
      {"id": "Q1", "type": "QCM", "points": 5, "estimated_time_minutes": 8},
      {"id": "Q2", "type": "Exercice d'application", "points": 10, "estimated_time_minutes": 20}
    ],
    "aa_mapping": [
      {"question_id": "Q1", "aa_ids": ["AA1"]},
      {"question_id": "Q2", "aa_ids": ["AA2"]}
    ],
    "bloom_mapping": [
      {"question_id": "Q1", "level": "Remember"},
      {"question_id": "Q2", "level": "Apply"}
    ],
    "duration_estimation": {
      "total_minutes": 95,
      "overload_flag": false
    },
    "quality_indicators": {
      "aa_coverage": true,
      "instruction_clarity": true,
      "marking_transparency": false
    },
    "improvement_proposals": [
      "Clarifier la consigne de Q2",
      "Ajouter une activite pratique couvrant AA3"
    ],
    "generated_items": [
      {"id": "NEW_EX1", "type": "Exercice", "points": 8},
      {"id": "NEW_PR1", "type": "Pratique", "points": 12}
    ],
    "reorganization_plan": {
      "remove": ["Q3"],
      "insert": ["NEW_EX1", "NEW_PR1"],
      "renumber": true
    },
    "official_report": {
      "template_name": "rapport examen officiel.tex",
      "output_tex": "rapport_evaluation_final.tex",
      "output_pdf": "rapport_evaluation_final.pdf"
    },
    "notifications": {
      "recipients": ["teacher1@esprit.tn", "teacher2@esprit.tn"],
      "status": "ready_to_send"
    }
  }
}
\end{lstlisting}

\section*{\colorbox{EspritBleu}{\parbox{\dimexpr\linewidth-2\fboxsep}{\centering \color{white}6. Structure des Dossiers Recommandée}}}

\begin{lstlisting}
exam-ai-pipeline/
├── inputs/
│   ├── exam.pdf
│   ├── template.tex
│   ├── rapport examen officiel.tex
│   ├── syllabus.json
│   ├── learning_outcomes.json
│   └── institution_rules.json
├── outputs/
│   ├── extracted_text.json
│   ├── questions.json
│   ├── classification.json
│   ├── marking_scheme.json
│   ├── aa_mapping.json
│   ├── bloom.json
│   ├── duration_estimation.json
│   ├── quality.json
│   ├── improvement_proposals.json
│   ├── generated_items.json
│   ├── reorganization_plan.json
│   ├── official_report.json
│   ├── notification_log.json
│   ├── platform_export.json
│   ├── consolidated_report.json
│   ├── report_final.tex
│   ├── rapport_evaluation_final.tex
│   ├── report_final.pdf
│   └── rapport_evaluation_final.pdf
├── agents/
│   ├── orchestrator.py
│   ├── ingestion_agent.py
│   ├── ocr_agent.py
│   ├── metadata_agent.py
│   ├── format_agent.py
│   ├── segmentation_agent.py
│   ├── classification_agent.py
│   ├── marking_agent.py
│   ├── aa_mapping_agent.py
│   ├── bloom_agent.py
│   ├── question_duration_agent.py
│   ├── quality_agent.py
│   ├── difficulty_agent.py
│   ├── improvement_agent.py
│   ├── generation_agent.py
│   ├── reorganization_agent.py
│   ├── compliance_agent.py
│   ├── consolidation_agent.py
│   ├── official_report_agent.py
│   ├── latex_agent.py
│   ├── validator_agent.py
│   ├── platform_exam_agent.py
│   └── teacher_notification_agent.py
└── prompts/
    ├── classification_prompt.txt
    ├── aa_mapping_prompt.txt
    ├── bloom_prompt.txt
    ├── quality_prompt.txt
    ├── improvement_prompt.txt
    ├── generation_prompt.txt
    ├── reorganization_prompt.txt
    ├── official_report_prompt.txt
    └── notification_prompt.txt
\end{lstlisting}

\section*{\colorbox{EspritBleu}{\parbox{\dimexpr\linewidth-2\fboxsep}{\centering \color{white}7. Workflow Opérationnel}}}

\begin{longtable}{|p{1.2cm}|p{5cm}|p{8cm}|}
\hline
\rowcolor{GrisClair}
\textbf{\#} & \textbf{Étape} & \textbf{Description} \\ \hline
1 & Ingestion & Chargement du fichier examen et des modèles \texttt{.tex} \\ \hline
2 & OCR / Reconstruction & Extraction ou reconstruction du texte \\ \hline
3 & Métadonnées & Extraction des informations générales \\ \hline
4 & Forme & Analyse des consignes pratiques de passation \\ \hline
5 & Segmentation & Découpage en questions, exercices et sous-parties \\ \hline
6 & Classification & Détermination du type de chaque élément \\ \hline
7 & Barème & Détection des points et du barème détaillé \\ \hline
8 & Mapping AA & Association questions / acquis d’apprentissage \\ \hline
9 & Bloom & Classification cognitive des questions \\ \hline
10 & Durée par question & Estimation du temps de traitement par élément \\ \hline
11 & Qualité & Calcul des indicateurs qualité \\ \hline
12 & Difficulté & Estimation de la difficulté \\ \hline
13 & Améliorations & Génération des propositions d’amélioration \\ \hline
14 & Génération & Création de nouvelles questions, exercices et pratiques \\ \hline
15 & Réorganisation & Suppression, insertion, remplacement, renumérotation \\ \hline
16 & Conformité & Vérification institutionnelle \\ \hline
17 & Consolidation & Fusion des résultats en JSON unique \\ \hline
18 & Rapport officiel & Génération via \texttt{rapport examen officiel.tex} \\ \hline
19 & Génération LaTeX & Injection dans \texttt{template.tex} \\ \hline
20 & Validation & Contrôle et correction des fichiers \texttt{.tex} \\ \hline
21 & Plateforme & Génération de l’examen sur la plateforme cible \\ \hline
22 & Notification & Envoi à tous les enseignants du cours \\ \hline
\end{longtable}

\section*{\colorbox{EspritBleu}{\parbox{\dimexpr\linewidth-2\fboxsep}{\centering \color{white}8. Exemple de Variables Injectées dans les Modèles .tex}}}

\subsection*{Pour le template de l’examen}

\begin{lstlisting}
\textbf{Nom de l'epreuve :} {{ exam_header.exam_name }}
\textbf{Classe :} {{ exam_header.class }}
\textbf{Date :} {{ exam_header.date }}
\textbf{Duree annoncee :} {{ exam_header.duration }}
\textbf{Duree estimee :} {{ duration_estimation.total_minutes }}
\end{lstlisting}

\subsection*{Pour le fichier \texttt{rapport examen officiel.tex}}

\begin{lstlisting}
\textbf{Nom de l'epreuve :} {{ exam_header.exam_name }}
\textbf{Enseignants :} {{ exam_header.teachers }}
\textbf{Qualite :} {{ quality_indicators }}
\textbf{Mapping AA :} {{ aa_mapping }}
\textbf{Bloom :} {{ bloom_mapping }}
\textbf{Ameliorations :} {{ improvement_proposals }}
\end{lstlisting}

\subsection*{Exemple de boucle pour les questions}

\begin{lstlisting}
{% for q in questions %}
{{ q.id }} & {{ q.type }} & {{ q.points }} & {{ q.bloom_level }} & {{ q.estimated_time_minutes }} \\
{% endfor %}
\end{lstlisting}

\section*{\colorbox{EspritBleu}{\parbox{\dimexpr\linewidth-2\fboxsep}{\centering \color{white}9. Recommandations de Mise en Production}}}

\begin{itemize}[leftmargin=1cm]
    \item Utiliser un pipeline en cinq couches :
    \begin{enumerate}
        \item extraction documentaire,
        \item analyse IA,
        \item amélioration et génération,
        \item consolidation JSON,
        \item génération LaTeX, publication et notification.
    \end{enumerate}
    \item Séparer les agents d’analyse, de rendu et de communication.
    \item Rendre chaque agent testable indépendamment.
    \item Journaliser toutes les décisions critiques.
    \item Prévoir un niveau de confiance par agent.
    \item Autoriser une validation humaine avant compilation finale, publication ou notification.
\end{itemize}

\section*{\colorbox{EspritBleu}{\parbox{\dimexpr\linewidth-2\fboxsep}{\centering \color{white}10. Conclusion}}}

L’architecture proposée permet de transformer un examen réel en :
\begin{itemize}[leftmargin=1cm]
    \item un rapport officiel d’évaluation,
    \item une nouvelle proposition d’examen enrichie et réorganisée,
    \item un export exploitable sur une plateforme cible,
    \item une notification automatisée vers tous les enseignants concernés.
\end{itemize}

Elle sépare clairement :
\begin{itemize}[leftmargin=1cm]
    \item l’analyse de l’existant,
    \item l’évaluation pédagogique,
    \item la génération d’améliorations,
    \item la production de contenus nouveaux,
    \item la génération du rapport officiel,
    \item la publication et la communication.
\end{itemize}

\vspace{0.8cm}

\begin{center}


% =========================
% TABLEAU RECAPITULATIF DES AGENTS
% =========================
\section*{\colorbox{EspritBleu}{\parbox{\dimexpr\linewidth-2\fboxsep}{\centering \color{white}Tableau Récapitulatif des Agents IA}}}

\renewcommand{\arraystretch}{1.4}
\small

\begin{longtable}{|p{3.2cm}|p{3.2cm}|p{3.2cm}|p{4.8cm}|p{2.8cm}|}
\hline
\rowcolor{GrisClair}
\textbf{Agent} & \textbf{Entrées} & \textbf{Sorties} & \textbf{Rôle} & \textbf{Dépendances} \\ \hline
\endfirsthead

\hline
\rowcolor{GrisClair}
\textbf{Agent} & \textbf{Entrées} & \textbf{Sorties} & \textbf{Rôle} & \textbf{Dépendances} \\ \hline
\endhead

Orchestrateur Central
& Examen, \texttt{template.tex}, \texttt{rapport examen officiel.tex}, AA, règles, métadonnées cours
& JSON consolidé, fichiers finaux, journal
& Pilote l’ensemble du workflow, ordonne les appels, gère erreurs et validations
& Tous les agents \\ \hline

Agent d’Ingestion Documentaire
& PDF, DOCX, image, texte
& Texte brut, pages, blocs détectés
& Détecte le format et extrait le contenu initial
& Aucun \\ \hline

Agent OCR / Reconstruction
& Document scanné, images, texte brut
& Texte reconstruit, structure logique
& Reconstruit les blocs du document si extraction directe insuffisante
& Ingestion \\ \hline

Agent d’Extraction de Métadonnées
& Texte reconstruit, en-têtes, consignes
& Nom épreuve, classe, langue, durée, date, enseignants, pages
& Extrait les informations générales de l’épreuve
& Ingestion, OCR \\ \hline

Agent d’Analyse de la Forme de l’Examen
& Consignes, règles institutionnelles
& Réponse sur feuille, documents autorisés, calculatrice, PC, Internet, validation
& Analyse les modalités pratiques de passation
& Métadonnées, OCR \\ \hline

Agent de Segmentation des Questions
& Texte structuré de l’examen
& Questions, sous-questions, exercices, identifiants
& Découpe l’examen en unités évaluatives
& OCR \\ \hline

Agent de Classification Automatique des Questions
& Questions segmentées
& Type de chaque question
& Classe les éléments en QCM, rédaction, pratique, étude de cas, etc.
& Segmentation \\ \hline

Agent d’Extraction du Barème
& Questions, texte de l’examen
& Points par question, total, barème détaillé ou non
& Détecte la répartition des points
& Segmentation \\ \hline

Agent de Mapping AA
& Questions, syllabus, acquis d’apprentissage
& Matrice Question $\rightarrow$ AA
& Associe chaque question aux acquis d’apprentissage
& Segmentation, Classification \\ \hline

Agent Bloom Taxonomy
& Questions, objectifs, verbes d’action
& Niveau Bloom par question, distribution globale
& Attribue un niveau cognitif à chaque question
& Segmentation, Classification \\ \hline

Agent d’Estimation de la Durée par Question
& Questions, difficulté, barème, niveau étudiants
& Temps estimé par question, temps total
& Estime la durée de traitement de chaque élément
& Segmentation, Classification, Barème \\ \hline

Agent d’Indicateurs Qualité
& AA, Bloom, barème, durée, types de questions
& Indicateurs qualité
& Calcule les indicateurs pédagogiques et qualitatifs
& Mapping AA, Bloom, Barème, Durée \\ \hline

Agent d’Analyse de Difficulté
& Questions, durée estimée, barème
& Difficulté par question, difficulté globale
& Évalue la complexité et le niveau global de l’examen
& Classification, Durée, Barème \\ \hline

Agent de Génération des Propositions d’Amélioration
& Indicateurs qualité, difficulté, durée, AA
& Liste des améliorations proposées
& Suggère corrections, reformulations, suppressions ou ajouts
& Qualité, Difficulté, Durée, Mapping AA \\ \hline

Agent de Génération de Questions, d’Exercices et de Pratiques
& Objectifs pédagogiques, AA, contraintes, améliorations
& Nouvelles questions, exercices, pratiques
& Génère de nouveaux contenus d’évaluation adaptés
& Améliorations, Mapping AA, Bloom \\ \hline

Agent de Réorganisation de la Nouvelle Proposition
& Ancien examen, nouveaux contenus, plan d’amélioration
& Nouvelle structure d’examen
& Supprime, remplace, insère et renumérote les éléments
& Génération, Segmentation, Améliorations \\ \hline

Agent de Conformité Institutionnelle
& Examen, rapport, règles institutionnelles
& Rapport de conformité, éléments manquants
& Vérifie la conformité académique et documentaire
& Forme, Métadonnées, Réorganisation \\ \hline

Agent de Consolidation des Résultats
& Sorties de tous les agents
& JSON unique consolidé
& Fusionne toutes les analyses dans une structure unique
& Tous les agents précédents \\ \hline

Agent Générateur du Rapport d’Évaluation d’Examen
& \texttt{rapport examen officiel.tex}, JSON consolidé
& \texttt{rapport\_evaluation\_final.tex}, PDF
& Produit le rapport officiel d’évaluation d’examen
& Consolidation \\ \hline

Agent de Génération LaTeX
& \texttt{template.tex}, JSON consolidé
& \texttt{report\_final.tex}
& Injecte les données dans le modèle LaTeX de sortie
& Consolidation \\ \hline

Agent de Validation LaTeX
& Fichiers \texttt{.tex} générés
& Fichiers corrigés, logs de validation
& Vérifie et corrige les erreurs de compilation LaTeX
& Générateur rapport, Génération LaTeX \\ \hline

Agent de Génération de l’Examen sur la Plateforme
& Nouvelle structure d’examen, métadonnées, barème
& Export plateforme, payload JSON/XML
& Génère l’examen dans le format attendu par la plateforme cible
& Réorganisation, Consolidation \\ \hline

Agent de Notification aux Enseignants du Cours
& Liste enseignants, liens, statuts, rapports
& Notifications envoyées, historique
& Informe tous les enseignants du cours des mises à jour et livrables
& Rapport officiel, Plateforme, Orchestrateur \\ \hline

\end{longtable}
\normalsize
\textbf{Fin du document}
\end{center}

\end{document}